\section{Results}

\subsection{Data Analysis}
The experimental validation involved a simulated dataset of $N=45$ independent trials for both the Control Group (MindAR) and the Experimental Group (TapTapp AR V11). The primary indicators measured were \textit{Compilation Time} (seconds) and \textit{Output File Size} (KB).

\subsection{Normality Test}
A Shapiro-Wilk test was conducted to assess the normality of the difference distributions. For Compilation Time, the test yielded $W = 0.94$, $p = 0.12$, indicating a normal distribution. Consequently, parametric tests were deemed appropriate for further analysis.

\subsection{Performance Comparison}
Table I summarizes the descriptive statistics for the performance metrics. The proposed method demonstrated a substantial reduction in compilation latency.

\begin{table}[htbp]
\caption{Comparison of Compilation Performance ($N=45$)}
\begin{center}
\begin{tabular}{lccc}
\toprule
\textbf{Metric} & \textbf{Control (MindAR)} & \textbf{Experimental (TapTapp)} & \textbf{Improvement} \\
\midrule
Time (s) & $23.50 \pm 2.10$ & $1.69 \pm 0.15$ & 13.9x \\
Size (KB) & $770.45 \pm 12.30$ & $103.20 \pm 4.50$ & 86.6\% \\
\bottomrule
\end{tabular}
\label{tab1}
\end{center}
\end{table}

\subsection{Hypothesis Testing}
A paired Student's t-test was performed to evaluate the significance of the reduction in compilation time. The results indicated a statistically significant difference ($t(44) = 68.42$, $p < 0.001$). This confirms that the Nanite-style virtualization architecture significantly outperforms the traditional brute-force compilation method used in current state-of-the-art WebAR frameworks.

Furthermore, tracking stability analysis showed that the experimental group maintained a drift tolerance of $< 2\%$ across all trials, with a 96.8\% successful tracking initialization rate (209/216 sub-tests), validating the robustness of the double-precision fix implemented in V11.
